%%%%%%%%%%%%%%%%%%%%%%%%%%%%%%%%%%%%%%%%%%%%%%%%%%%%%%%%%%%%%%%%%%%%%%%%%%% 
% 
% Generic template for TFC/TFM/TFG/Tesis
% 
% By:
%  + Javier Macías-Guarasa.
%    Departamento de Electrónica
%    Universidad de Alcalá
%  + Roberto Barra-Chicote.
%    Departamento de Ingeniería Electrónica
%    Universidad Politécnica de Madrid
% 
% By: + Javier Macías-Guarasa. Departamento de Electrónica Universidad de Alcalá + Roberto Barra-Chicote. Departamento de Ingeniería Electrónica Universidad Politécnica de Madrid
% 
% Based on original sources by Roberto Barra, Manuel Ocaña, Jesús Nuevo, Pedro Revenga, Fernando Herránz and Noelia Hernández. Thanks a lot to all of them, and to the many anonymous contributors found (thanks to google) that provided help in setting all this up.
% 
% See also the additionalContributors.txt file to check the name of additional contributors to this work.
% 
% If you think you can add pieces of relevant/useful examples, improvements, please contact us at (macias@depeca.uah.es)
% 
% You can freely use this template and please contribute with comments or suggestions!!!
% 
%%%%%%%%%%%%%%%%%%%%%%%%%%%%%%%%%%%%%%%%%%%%%%%%%%%%%%%%%%%%%%%%%%%%%%%%%%% 

\chapter{Introducción}
\label{cha:introduccion}

\begin{FraseCelebre}
  \begin{Frase}
    Desocupado lector, sin juramento me podrás creer que quisiera que este libro [...] fuera el más hermoso, el más gallardo y más discreto que pudiera imaginarse\footnote{Tomado de ejemplos del proyecto \texis{}.}.
  \end{Frase}
  \begin{Fuente}
    Miguel de Cervantes, Don Quijote de la Mancha
  \end{Fuente}
\end{FraseCelebre}


\section{Presentación}
\label{sec:presentacion}

Esta plantilla proporciona una estructura común para redactar trabajos académicos con \LaTeX{}. Incluye soporte para trabajos fin de carrera de planes anteriores, trabajos fin de grado, trabajos fin de máster y tesis doctorales. La mayor parte de las portadas y documentos administrativos están adaptados a titulaciones de la Universidad de Alcalá, aunque también se conserva soporte para determinados programas de doctorado de la Universidad Politécnica de Madrid.

El repositorio contiene, además del libro principal, plantillas para el anteproyecto y para distintos documentos administrativos asociados a la presentación, depósito, evaluación y publicación de los trabajos. La configuración compartida permite reutilizar los datos del trabajo, del autor y de los tutores sin tener que introducirlos de nuevo en cada documento.

La plantilla pretende facilitar tanto la composición del documento como el cumplimiento de los formatos institucionales. No obstante, estos formatos pueden cambiar con el tiempo, por lo que conviene comprobar siempre la normativa vigente de la titulación antes de entregar la documentación.

\section{Organización del repositorio}
\label{sec:organizacion-repositorio}

Los principales directorios del repositorio son:

\begin{description}
  \item[\texttt{Book/}] Documento académico principal, incluidas las cubiertas, capítulos, anexos, bibliografía, resúmenes y demás elementos del libro.
  \item[\texttt{Config/}] Configuración compartida, selección de la titulación y datos personales y académicos definidos en \texttt{myconfig.tex}.
  \item[\texttt{Anteproyecto/}] Plantilla para preparar la propuesta o anteproyecto del trabajo.
  \item[\texttt{PapeleoTFG/}] Formularios administrativos relacionados con los trabajos fin de grado.
  \item[\texttt{PapeleoTFM/}] Formularios administrativos relacionados con los trabajos fin de máster.
  \item[\texttt{PapeleoPHD/}] Formularios administrativos relacionados con las tesis doctorales.
  \item[\texttt{normativas/}] Normativas y documentos institucionales de referencia.
  \item[\texttt{UsefulDocs/}] Documentación auxiliar útil para preparar y revisar el trabajo.
  \item[\texttt{AdminScripts/}] Herramientas de mantenimiento utilizadas para administrar la plantilla.
  \item[\texttt{Deprecated/}] Material antiguo conservado únicamente como referencia o por compatibilidad.
\end{description}

La configuración general se realiza en \texttt{Config/myconfig.tex}. Los distintos documentos del repositorio leen este archivo para mantener coherentes los nombres, títulos, fechas y demás datos compartidos.



\section{Uso de la plantilla}
\label{sec:uso-generico-de}


\subsection{Prerrequisitos}
\label{sec:prerrequisitos}

Para compilar la plantilla se necesita una distribución de \LaTeX{} suficientemente completa y actualizada. TeX Live es una opción habitual en GNU/Linux, MiKTeX y TeX Live están disponibles para Windows, y MacTeX proporciona una distribución adaptada a macOS. La distribución elegida debe proporcionar \texttt{pdflatex} y Biber, además de los paquetes cargados desde \texttt{Config/preamble.tex}.

Biber es necesario cuando el documento incluye bibliografía. Si se utilizan acrónimos o símbolos, también se necesita \texttt{makeglossaries}; el libro de ejemplo distribuido con la plantilla contiene ambos elementos y, por tanto, requiere esta herramienta para generar todos sus listados.

Si durante la compilación se informa de que falta un fichero \texttt{.sty}, se debe instalar el paquete de la distribución de \LaTeX{} que lo contiene. En MiKTeX puede activarse la instalación automática de paquetes; TeX Live y MacTeX proporcionan sus propias herramientas para administrar los paquetes instalados.

La plantilla no requiere el uso de \texttt{make}. Puede compilarse desde la terminal o con cualquier editor o herramienta de construcción de \LaTeX{} que permita ejecutar la secuencia descrita en la sección~\ref{sec:compilacion}. El \texttt{Makefile} incluido es únicamente una forma opcional de automatizar el proceso y resulta especialmente cómodo en sistemas GNU/Linux.

\subsection{Compilación}
\label{sec:compilacion}

Antes de compilar el documento debe completarse la configuración compartida en \texttt{Config/myconfig.tex}. La secuencia completa utilizada por el libro de ejemplo es:

\begin{verbatim}
pdflatex book.tex
biber book
makeglossaries book
pdflatex book.tex
pdflatex book.tex
\end{verbatim}

Biber genera la bibliografía y \texttt{makeglossaries} genera los listados de acrónimos y símbolos. Si un documento personalizado no contiene bibliografía, acrónimos ni símbolos, se pueden omitir las herramientas correspondientes. Las pasadas adicionales de \texttt{pdflatex} permiten actualizar correctamente las referencias cruzadas, los índices, la bibliografía y, cuando existan, los glosarios.

\subsubsection{Compilación desde un editor de \LaTeX{}}

Se puede utilizar TeXstudio, TeXShop u otro editor de \LaTeX{} en cualquiera de los sistemas operativos compatibles. El editor debe configurarse para utilizar \texttt{pdflatex} como compilador, Biber como procesador de la bibliografía y, únicamente si se utilizan acrónimos o símbolos, \texttt{makeglossaries} para generar sus listados. Después deben realizarse las dos pasadas finales de \texttt{pdflatex}.

Algunos editores permiten definir una orden o cadena de construcción personalizada que ejecuta automáticamente todo el proceso. Otros proporcionan órdenes separadas para \LaTeX{}, la bibliografía y los glosarios; en ese caso deben ejecutarse respetando el orden indicado anteriormente. El fichero principal que debe abrirse y compilarse es \texttt{Book/book.tex}.

\subsubsection{GNU/Linux}

En GNU/Linux se puede compilar el documento con el editor de \LaTeX{} preferido o ejecutar directamente desde una terminal la secuencia indicada anteriormente. Los comandos deben ejecutarse desde el directorio \texttt{Book/}.

Si se dispone de la utilidad \texttt{make}, puede automatizarse todo el proceso mediante:

\begin{verbatim}
cd Book
make
\end{verbatim}

El uso de \texttt{make} es opcional y no modifica los requisitos propios de \LaTeX{}: se limita a ejecutar automáticamente las herramientas necesarias y a generar las copias finales.

\subsubsection{Windows}

En Windows se puede utilizar MiKTeX o TeX Live junto con TeXstudio u otro editor de \LaTeX{}. En la configuración del editor se debe seleccionar Biber como procesador de la bibliografía y, si se utilizan acrónimos o símbolos, añadir \texttt{makeglossaries} a la cadena de construcción. La ejecución de \texttt{makeglossaries} también requiere disponer de un intérprete de Perl; Strawberry Perl es una de las alternativas disponibles para Windows.

No es necesario instalar \texttt{make} para trabajar con la plantilla en Windows. Si se desea utilizar el \texttt{Makefile}, puede hacerse desde un entorno compatible que proporcione las herramientas requeridas, como WSL, pero normalmente resulta más sencillo configurar la cadena de compilación en el propio editor.

Parte de las instrucciones originales para configurar TeXmaker en Windows fueron aportadas por Cristina Losada.

\subsubsection{macOS}

En macOS se puede utilizar MacTeX junto con TeXShop, TeXstudio u otro editor compatible. La cadena de construcción debe incluir \texttt{pdflatex}, Biber y, si se utilizan acrónimos o símbolos, \texttt{makeglossaries} en el orden indicado anteriormente.

Si el editor no actualiza correctamente los acrónimos o símbolos, se puede ejecutar manualmente \texttt{makeglossaries book} desde el directorio \texttt{Book/} y realizar después dos nuevas pasadas de \texttt{pdflatex}. No es necesario borrar el resto de los ficheros del proyecto para regenerar los glosarios.

Las primeras comprobaciones de la generación de glosarios en macOS fueron aportadas por Carlos Cruz.

\subsubsection{Compilación mediante el \texttt{Makefile}}

El objetivo predeterminado \texttt{all}, ejecutado al invocar \texttt{make} sin argumentos, convierte primero los diagramas y figuras que lo necesiten. A continuación ejecuta \texttt{pdflatex}, Biber, \texttt{makeglossaries} y las dos pasadas finales de \texttt{pdflatex}. Dado que el libro de ejemplo utiliza acrónimos y símbolos, el \texttt{Makefile} ejecuta siempre \texttt{makeglossaries}.

La compilación genera \texttt{book.pdf} y \texttt{book-compressed.pdf}. También crea copias cuyos nombres identifican el tipo de trabajo, la titulación, el autor y el idioma configurados. La versión comprimida utiliza la configuración \texttt{/prepress} de Ghostscript; la antigua versión de calidad \texttt{/screen} está desactivada porque degradaba excesivamente las imágenes.

El objetivo \texttt{all} también actualiza el fichero \texttt{RELEASE.txt} con la etiqueta más reciente del repositorio e intenta añadirlo al índice de Git. Por tanto, después de compilar conviene revisar el estado del repositorio antes de preparar un commit.

Los principales objetivos disponibles en el \texttt{Makefile} son:

\begin{description}
  \item[\texttt{all}] Compila el documento completo y genera las copias finales; es el objetivo ejecutado de forma predeterminada por \texttt{make}.
  \item[\texttt{clean}] Elimina los ficheros auxiliares y los PDF de trabajo generados durante la compilación.
  \item[\texttt{all\_latexmk}] Ofrece un procedimiento alternativo y experimental de compilación basado en \texttt{latexmk}.
  \item[\texttt{clean\_latexmk}] Elimina los ficheros auxiliares mediante \texttt{latexmk} y completa después la limpieza de otros ficheros generados.
  \item[\texttt{snapshot}] Genera mediante \texttt{latexpand} una copia unificada del documento que puede utilizarse como referencia para una revisión posterior.
  \item[\texttt{flatten}] Genera una copia unificada del estado actual del documento.
  \item[\texttt{latexdiff}] Compara la instantánea con el estado actual y genera un PDF que muestra las diferencias.
  \item[\texttt{backup-chapters}] Crea copias fechadas de los capítulos y apéndices.
  \item[\texttt{orig-chapters}] Sustituye los capítulos y apéndices de trabajo por las versiones completas de ejemplo, después de crear una copia de seguridad.
  \item[\texttt{bare-chapters}] Sustituye los capítulos y apéndices de trabajo por versiones mínimas, después de crear una copia de seguridad.
  \item[\texttt{bare}] Prepara una estructura mínima mediante \texttt{bare-chapters} y ofrece eliminar los gráficos y diagramas existentes.
  \item[\texttt{tar}] Genera archivos comprimidos con los fuentes y recursos de la plantilla.
\end{description}

Los objetivos \texttt{bare}, \texttt{bare-chapters} y \texttt{orig-chapters} modifican el contenido de los directorios de capítulos o apéndices. Además, \texttt{bare} puede eliminar los ficheros de \texttt{figures/} y \texttt{diagrams/} si se confirma expresamente durante su ejecución. Deben utilizarse únicamente al preparar una copia nueva del proyecto y después de comprobar que los cambios importantes están guardados en el sistema de control de versiones.

\subsubsection{Herramientas y prestaciones opcionales}

Además de la compilación básica, el repositorio permite utilizar varias herramientas opcionales:

\begin{description}
  \item[Ghostscript] Genera la copia comprimida del documento cuando se utiliza el objetivo \texttt{all} del \texttt{Makefile}.
  \item[Inkscape] Convierte automáticamente los diagramas en formato \texttt{svg} a \texttt{pdf}.
  \item[Dia y \texttt{epspdf}] Convierten automáticamente los diagramas en formato \texttt{dia} a \texttt{pdf}.
  \item[\texttt{epspdf}] Convierte automáticamente las figuras en formato \texttt{eps} a \texttt{pdf}.
  \item[\texttt{latexmk}] Proporciona el procedimiento alternativo utilizado por los objetivos \texttt{all\_latexmk} y \texttt{clean\_latexmk}.
  \item[\texttt{latexpand} y \texttt{latexdiff}] Permiten unificar los fuentes y generar documentos que muestran las diferencias entre dos versiones.
  \item[Perl] Es necesario para ejecutar \texttt{makeglossaries} en los sistemas que no incluyan ya un intérprete compatible.
\end{description}

Estas herramientas solo deben instalarse cuando se vaya a utilizar la prestación correspondiente. No son requisitos para redactar ni compilar un documento básico que no dependa de ellas.

\subsection{Estructura del documento generado por la plantilla}
\label{sec:estr-del-docum}

La estructura del documento principal se define en \texttt{Book/book.tex}. En este fichero se pueden añadir, eliminar, comentar o reordenar los distintos elementos mediante sus correspondientes órdenes \texttt{\textbackslash{}input} o \texttt{\textbackslash{}includepdf}. Los ejemplos incluidos no constituyen una estructura obligatoria y deben adaptarse a la normativa y a las necesidades de cada trabajo.

El documento se divide en tres partes gestionadas por las órdenes \texttt{\textbackslash{}frontmatter}, \texttt{\textbackslash{}mainmatter} y \texttt{\textbackslash{}backmatter}:

\begin{description}
  \item[\texttt{\textbackslash{}frontmatter}] Contiene las cubiertas, los documentos preliminares, la dedicatoria, los agradecimientos, los resúmenes y los distintos índices y listados.
  \item[\texttt{\textbackslash{}mainmatter}] Contiene los capítulos principales, los apartados técnicos opcionales, la bibliografía y los apéndices.
  \item[\texttt{\textbackslash{}backmatter}] Contiene la contraportada cuando corresponde al tipo de trabajo seleccionado.
\end{description}

La estructura distribuida con la plantilla incluye los siguientes elementos:

\begin{itemize}
  \item Portada y páginas iniciales, seleccionadas automáticamente por \texttt{cover/cover.tex} a partir de la titulación configurada mediante \texttt{\textbackslash{}myDegree}. Las distintas plantillas institucionales se encuentran en el directorio \texttt{cover/}.
  \item Documentos PDF preliminares, que pueden insertarse con \texttt{\textbackslash{}includepdf}. El fichero \texttt{letters/sampleLetter.pdf} muestra cómo incorporar todas las páginas de un documento, y \texttt{letters/sampleLetter-pages.pdf} permite probar la selección de páginas concretas. Estos documentos pueden ser necesarios, por ejemplo, en determinadas tesis doctorales o trabajos fin de máster.
  \item Dedicatoria, incluida desde \texttt{dedication/dedicatoria.tex}. Puede editarse o comentarse su orden \texttt{\textbackslash{}input} si no se desea utilizar.
  \item Agradecimientos, incluidos desde \texttt{acknowledgements/agradecimientos.tex}. También pueden editarse u omitirse.
  \item Definiciones de acrónimos, incluidas desde \texttt{acronyms/defacronymsgl.tex}, y listado correspondiente, configurado en \texttt{acronyms/acronymsgl.tex}. Ambos elementos pueden comentarse cuando no se utilicen acrónimos.
  \item Definiciones de símbolos, incluidas desde \texttt{symbols/defsymbolsgl.tex}, y listado correspondiente, configurado en \texttt{symbols/symbolsgl.tex}. Ambos elementos pueden comentarse cuando no se utilicen símbolos.
  \item Resumen en español, incluido desde \texttt{abstract/resumen.tex}, y resumen en inglés, incluido desde \texttt{abstract/abstract.tex}.
  \item Resumen extendido en español, disponible en \texttt{abstract/resumen-extendido.tex}. Su inclusión está comentada de forma predeterminada y puede activarse cuando lo requiera el tipo de trabajo.
  \item Índice general, índice de figuras e índice de tablas, generados desde \texttt{cover/toc+lof+lot.tex}.
  \item Listados adicionales, configurados en \texttt{cover/extralistings.tex}. El ejemplo incluye listados para fragmentos de código, algoritmos y vídeos, y puede editarse o comentarse si no se necesita.
  \item Capítulos principales, incluidos desde el directorio \texttt{chapters/}. El ejemplo proporciona una estructura formada por introducción, estudio teórico, desarrollo, resultados y conclusiones, que debe adaptarse libremente al contenido real del trabajo.
  \item Pliego de condiciones, disponible en \texttt{pliego/pliego.tex}, y presupuesto, disponible en \texttt{presupuesto/presupuesto.tex}. Sus órdenes de inclusión están comentadas de forma predeterminada y deben activarse únicamente cuando sean necesarios. Los ejemplos proceden de un trabajo fin de carrera de Jesús Martínez y se conservan como material de referencia; debido a su antigüedad, deben contrastarse con la normativa y los formatos actuales.
  \item Bibliografía, impresa mediante \texttt{biblio/bibliography.tex}. Los ficheros \texttt{.bib} que se procesarán se declaran en \texttt{biblio/bibliofiles.tex} mediante las variables \texttt{\textbackslash{}mybibfileOne}, \texttt{\textbackslash{}mybibfileTwo} y las sucesivas variables disponibles.
  \item Apéndices, agrupados mediante el entorno \texttt{appendices}. El libro de ejemplo incluye \texttt{appendix/manual.tex}, \texttt{appendix/herramientas.tex} y \texttt{appendix/versiones.tex}.
  \item Contraportada, seleccionada automáticamente mediante \texttt{cover/backpage.tex} cuando existe una plantilla aplicable al tipo de trabajo y a la titulación configurados.
\end{itemize}

Además de editar los ficheros incluidos, se debe revisar \texttt{Book/book.tex} para comentar los elementos que no se necesiten y añadir los capítulos, apéndices o documentos externos propios. Cuando se incorpore un PDF mediante \texttt{\textbackslash{}includepdf}, conviene mantener la orden \texttt{\textbackslash{}clearemptydoublepage} posterior para que el siguiente elemento comience en la página adecuada.

\subsection{Definiciones específicas del tipo de documento}
\label{sec:definicion-del-tipo}

Para comenzar a usar la plantilla es fundamental revisar el fichero \texttt{book.tex} en el que se incluyen todos los detalles genéricos de la estructura usada en el documento, con comentarios que esperamos que os ayuden a entenderlo. Si sois de los impacientes, basta con que comencéis por la parte en la que se incluyen los distintos capítulos (buscad la parte de los \texttt{\textbackslash{}input\{chapters/*.tex\}}).

El fichero de configuración más importantes es el \texttt{Config/myconfig.tex} en el que se incluyen elementos para determinar la configuración específica de tu documento y que es de donde se toman los datos para las muchas automatizaciones del proceso de generación. El primero de ellos (para facilitar la generación de documentos en español o inglés), es el que define el idioma que vas a utilizar. Para seleccionarlo, basta con asignar \texttt{spanish} o \texttt{english} a la variable \texttt{\textbackslash{}myLanguage}. A partir de ella, el sistema generará las cabeceras y títulos adecuados a cada una.

Igualmente, tendréis que definir las siguientes variables sobre el tipo de trabajo y el autor\footnote{El fichero  \texttt{Config/myconfig.tex} ``sufre'' modificaciones frecuentes con la adición de nuevas variables para soportar un mayor número de automatismos (por ejemplo los que tienen que ver con los documento de papeleo necesarios), con lo que es seguro que encontrarás más variables en la versión actual del fichero que los que aquí hemos descrito. En cualquier caso, relax, porque los nombres son autoexplicativos.}:

\begin{itemize}
  \item Acrónimo de la titulación correspondiente al trabajo (variable \texttt{\textbackslash{}myDegree}): Que se seleccionará entre los definidos (por ahora\footnote{Este apartado se generó en marzo de 2020.} son \texttt{IT}, \texttt{IE}, \texttt{ITTSE}, \texttt{ITTST}, \texttt{ITI}, \texttt{GIEC}, \texttt{GIST}, \texttt{GITT}, \texttt{GIT}, \texttt{GIEAI}, \texttt{GITI}, \texttt{GIC}, \texttt{GII}, \texttt{GSI}, \texttt{GISI}, \texttt{MUSEA}, \texttt{MUIT}, \texttt{MUII}, \texttt{PHDUAH} y \texttt{PHDUPM}) y que automáticamente configura portadas\footnote{Un comentario sobre el color de las bandas en la portada de los TFGs: de acuerdo con la normativa, dicho color debe ser PANTONE 160c. Yo he intentado utilizar dicho color, pero el aspecto con el que finalmente aparece no es ni de lejos similar al del modelo que proporciona la EPS, con lo que he optado por utilizar el que se ve realmente en dicho modelo. Si quieres cambiarlo, busca ``PANTONE'' en \texttt{Config/preamble.tex}.}, entre otras cosas.

  \item Flag de división de directores (variable \texttt{\textbackslash{}myFlagSplittedAdvisors}). Si puesta a \texttt{false} indicará ``Tutores/Advisors'' en las páginas de portada. En caso contrario lo separará en líneas para ``Tutor/Advisor'' y ``Cotutor/Co-Advisor''.

  \item Especialidad (nuevo en TFGs desde el curso 2015) (variable \texttt{\textbackslash{}mySpecialty})

\end{itemize}

Sobre el documento en general:

\begin{itemize}
  \item Título del documento en español (variable \texttt{\textbackslash{}myBookTitleSpanish}).
  \item Título del documento en inglés (variable \texttt{\textbackslash{}myBookTitleEnglish}).
  \item Palabras clave en inglés (variable \texttt{\textbackslash{}myThesisKeywordsEnglish}).
  \item Palabras clave en español (variable \texttt{\textbackslash{}myThesisKeywords}).


\end{itemize}

Sobre el autor y el tutor y cotutor:

\begin{itemize}
  
  \item Nombre del autor del trabajo (variable \texttt{\textbackslash{}myAuthorName}).
  \item Apellidos del autor del trabajo (variable \texttt{\textbackslash{}myAuthorSurname}). La unión de esta definición con la del nombre formará automáticamente la variable \texttt{\textbackslash{}myAuthorFullName}
  \item Género del autor (variable \texttt{\textbackslash{}myAuthorGender}), ``male'' or ``female''.
  \item Email del autor (variable \texttt{\textbackslash{}myAuthorEmail}).
  \item DNI del autor del trabajo, usado en el papeleo de los TFGs (variable \texttt{\textbackslash{}mybookDNI}).

  \item También se ha iniciado soporte básico para generar la hoja de control de anteproyecto que se usa en la solicitud, con lo que hay variables relacionadas con los datos personales del alumno:

        \begin{itemize}
          \item Dirección: calle (variable \texttt{\textbackslash{}myAuthorStreet}), ciudad (variable \texttt{\textbackslash{}myAuthorCity}), código postal (variable \texttt{\textbackslash{}myAuthorPostalCode}), provincia (variable \texttt{\textbackslash{}myAuthorProvince}) y teléfono (variable \texttt{\textbackslash{}myAuthorTelephone}).
        \end{itemize}

  \item Nombre del primer (o único, en su caso) director/tutor del trabajo (variable \texttt{\textbackslash{}myAcademicTutorFullName}).
  \item Género del primer director/tutor (variable \texttt{\textbackslash{}myAcademicTutorGender}), ``male'' or ``female''.
  \item DNI del primer direector, para los permisos de publicación en abierto (variable \texttt{\textbackslash{}myFirstAdvisorDNI}

  
  \item Nombre del segundo (en su caso) director/tutor del trabajo (variable \texttt{\textbackslash{}myCoTutorFullName}).
  \item Género del segundo director/tutor (variable \texttt{\textbackslash{}myCoTutorGender}), ``male'' or ``female''.
  \item DNI del segundo director, para los permisos de publicación en abierto (variable \texttt{\textbackslash{}mySecondAdvisorDNI}), en su caso

\end{itemize}

Sobre la filiación del autor y el trabajo:

\begin{itemize}

  \item Centro en el que se realiza el trabajo (variable \texttt{\textbackslash{}mySchool}, que debería ser la Escuela Politénica Superior, pero se incluye por generalidad).
  \item Universidad en el que se realiza el trabajo (variable \texttt{\textbackslash{}myUniversity}, que debería ser la de Alcalá, pero se incluye por generalidad).
  \item Acrónimo de la Universidad (variable \texttt{\textbackslash{}myUniversityAcronym}).

  \item Departamento en el que se realiza el trabajo (variable \texttt{\textbackslash{}myDepartment}, en español, y \texttt{\textbackslash{}myDepartmentEnglish}, en inglés).
  
  \item Programa de Doctorado (en su caso) en el que se realiza el trabajo (variable \texttt{\textbackslash{}myPhDProgram}, en español, y \texttt{\textbackslash{}myPhDProgramEnglish}, en inglés).
  \item Grupo de investigación en el que se realiza el trabajo (variable \texttt{\textbackslash{}myResearchGroup}.

\end{itemize}

Sobre el tribunal:

\begin{itemize}

  \item Nombre del presidente del tribunal (variable \texttt{\textbackslash{}myTribunalPresident}).
  \item Nombre del primer vocal del tribunal (variable \texttt{\textbackslash{}myTribunalFirstSpokesperson}).
  \item Nombre del segundo vocal del tribunal (variable \texttt{\textbackslash{}myTribunalSecondSpokesperson}).
  \item Nombre del suplente del tribunal (variable \texttt{\textbackslash{}myTribunalAlternateMember}).
  \item Nombre del secretario del tribunal, en su caso (variable \texttt{\textbackslash{}myTribunalSecretary}).
  \item Nombre del Secretario del Departamento (que firmará parte del papeleo) (variable \texttt{\textbackslash{}myDepartmentSecretary}).

\end{itemize}

Sobre fechas varias:

\begin{itemize}
  
  \item Fecha del anteproyecto, en su caso (variable \texttt{\textbackslash{}myThesisProposalDate}), en el caso de que se use la plantilla del anteproyecto
  \item Fecha del anteproyecto en inglés, en su caso (variable \texttt{\textbackslash{}myThesisProposalDateEnglish}), en el caso de que se usen la plantilla del anteproyecto.

  \item Fecha del depósito del trabajo (variable \texttt{\textbackslash{}myThesisDepositDate})
  \item Fecha del depósito del trabajo en inglés (variable \texttt{\textbackslash{}myThesisDepositDateEnglish})

  \item Año del trabajo (variable \texttt{\textbackslash{}myThesisDefenseYear}).
  \item Fecha de la defensa del trabajo, en su caso (variable \texttt{\textbackslash{}myThesisDefenseDate}, en español, o \texttt{\textbackslash{}myThesisdefenseDateEnglish} en inglés).

  \item Fecha que aparecerá en la firma del papeleo (variable \texttt{\textbackslash{}myPaperworkDate}).
  
  \item Nombre del vicerrector (variable \texttt{\textbackslash{}myResearchVicerrector}).
  \item Género del vicerrector (variable \texttt{\textbackslash{}myResearchVicerrectorGender}).
\end{itemize}

Se ha comenzado a trabajar en la versión alfa del soporte para \textit{research reports}, y en ese caso se usa la variable \texttt{\textbackslash{}myResearchReportID}.

Y también hay otras variables de menor interés o en desuso en las que no vamos a entrar.

Parte de esa información se utilizará para rellenar la metainformación incluida en el fichero \texttt{pdf} que se genera.


También se definen los colores que se usarán en los hiperenlaces del documento:

\begin{itemize}
  \item Color de los enlaces en el índice de contenidos (variable \texttt{\textbackslash{}mytoclinkcolor}).
  \item Color de los enlaces en el índice de figuras (variable \texttt{\textbackslash{}myloflinkcolor}).
  \item Color de los enlaces en el índice de tablas (variable \texttt{\textbackslash{}mylotlinkcolor}).
  \item Color de los enlaces en otros índices (variable \texttt{\textbackslash{}myothertoclinkcolor}), de los que ahora se incluye un ejemplo en el fichero \texttt{cover/extralistings.tex}.
  \item Color de los enlaces (\texttt{\textbackslash{}ref}) en el documento (variable \texttt{\textbackslash{}mylinkcolor}).
  \item Color de los enlaces a URLs (variable \texttt{\textbackslash{}myurlcolor}).
  \item Color de los enlaces a referencias bibliográficas (variable \texttt{\textbackslash{}mycitecolor}).
\end{itemize}

Basta con que definas las variables correspondientes y la plantilla generará automáticamente las portadas adecuadas a la normativa y usará las definiciones específicas que hayas seleccionado.

Por si os hace falta, en \texttt{Config/postamble.tex} se definen muchas variables relacionadas con el tipo de trabajo, la personalización y adaptación a género y número en distintos documentos, etc. Por ejemplo las variables \texttt{\textbackslash{}myDegreefull} (igual a ``\myDegreefull'' en esta compilación), \texttt{\textbackslash{}myWorkType} (igual a ``\myWorkType'' en esta compilación) y \texttt{\textbackslash{}myWorkTypeFull} (igual a ``\myWorkTypeFull'' en esta compilación). Otro ejemplo sería el autor de contacto: \contactauthor.

Importante para las tesis de UAH: si necesitáis incluir ficheros pdf (los de autorización e informes de los tutores, por ejemplo), esta plantilla lo permite: mirad los \texttt{\textbackslash{}includepdf} en el \texttt{book.tex}.

\subsection{Plantilla de anteproyecto}
\label{sec:plantilla-de-anteproyecto}

Para el caso de los TFMs/TFGs/TFCs, se incluye una plantilla para realizar el anteproyecto.

La plantilla se encuentra en el directorio \texttt{Anteproyecto}, y en fichero \texttt{anteproyecto.tex} tenéis un ejemplo. El \texttt{Makefile} genera automáticamente el \texttt{anteproyecto.pdf}, haciendo un \texttt{make} en ese directorio, y tiene targets similares a los del \texttt{Makefile} del documento principal, incluyendo \texttt{flatten} y \texttt{latexdiff}, con lo que también puede generarse el fichero con control de cambios, tal y como se describe en la sección~\ref{sec:control-de-cambios} (cambiando las referencias a \texttt{book...} por \texttt{anteproyecto...}).



\subsection{Papeleo adicional}
\label{sec:introapp1}

Para el caso de los TFGs y TFMs se incluyen en los directorios \texttt{PapeleoTFG/} y \texttt{PapeleoTFM} algunos de los documentos que hay que generar en el momento de la defensa. En concreto, para TFGs:

\begin{itemize}
  \item La autorización del tutor para la publicación en abierto, en el fichero \texttt{PapeleoTFG/TFG-AutorizacionTutorPubAbierto.tex} para TFGs.
  \item La autorización del autor para la publicación en abierto, en el fichero \texttt{PapeleoTFG/TFG-AutorizacionAutorPubAbierto.tex}
  \item La autorización del tutor extranjero para la publicación en abierto, en el fichero \texttt{PapeleoTFG/TFG-AutorizacionTutorExtranjeroPubAbierto.tex}
  \item El informe y rúbrica del tutor para la defensa del TFG \texttt{PapeleoTFG/TFG-InformeTutorYRubrica.tex}
  \item El acta y rúbrica del tribunal para la defensa del TFG \texttt{PapeleoTFG/TFG-ActaYRubricaDefensa.tex}
\end{itemize}

Y para TFMs:

\begin{itemize}
  \item La autorización del tutor y autor para la publicación en abierto, en el fichero \texttt{PapeleoTFM/TFM-autorizacionPublicarAbierto.tex} para TFMs.
\end{itemize}

Para compilarlos, encontrarás un \texttt{Makefile} que genera los \texttt{pdfs} correspondientes a esos documentos. El sistema automáticamente adapta los términos masculinos/femeninos correspondientes y los singulares/plurales necesarios para el caso de que haya uno o varios directores/tutores.

Para el caso de las tesis doctorales, no está todavía preparado.


\subsection{Generación del documento con control de cambios (para revisión)}
\label{sec:control-de-cambios}

Uno de los problemas de \LaTeX{} frente a otros sistemas de edición de textos viene en el momento de la revisión de cambios realizados a un documento. En el caso que nos ocupa serían los que nos sugiere nuestro tutor de TFC/TFG/TFG o director de tesis doctoral y que luego querrá verificar cómo se han aplicado. En un procesador al estilo de libreoffice (vaaaaaaaaale, o Microsoft Word también), tenemos la opción de comparar documentos o llevar el control de cambios y que el sistema automáticamente nos marque los añadidos o borrados.

En \LaTeX{} también es posible si usamos un sistema de control de versiones (y si no lo estáis usando, deberíais plantearos seriamente el hacerlo), con la ayuda de herramientas adicionales.

Hay varias soluciones disponibles, la mayoría basada en el uso de la \texttt{latexdiff}~\cite{latexdiff} (o derivados). Lo recomendable sería el uso de un sistema automatizado como el disponible en \texttt{git-latexdiff} \cite{git-latexdiff}, pero necesita el soporte de \texttt{git}, y dados los orígenes en los que mantenía esto en \texttt{cvs} (flames to \texttt{/dev/null}, please), o pensando en los que no usaréis ningún sistema de control de versiones (si es así, mal hecho, por cierto, deberías usar alguno) el mecanismo que uso es un poco más rupestre (si quieres ayudarme a generalizar esto, sería perfecto).

\texttt{latexdiff} permite comparar dos versiones de un documento y generar un nuevo fichero fuente en \LaTeX{} que, al compilarlo, muestra un pdf ``bonito'', con las marcas de las diferencias (resaltando lo añadido y lo quitado). El resultado es perfecto para hacer una buena revisión de los cambios, y por supuesto no tiene punto de comparación con ver un \texttt{diff} a palo seco de las dos versiones del documento.

Lo que finalmente he implementado es muy ad-hoc, pero funciona\footnote{por algún motivo, a fecha de julio de 2021, \texttt{latexdiff} genera código erróneo, y he tenido que hacer una super ñapa en el \texttt{Makefile} para solucionarlo, así que en el futuro seguro que da problemas. Be warned!}.

El procedimiento para generar un documento que muestre las diferencias es el siguiente (asumo disponible un intérprete de comandos y la utilidad \texttt{make}):

\begin{itemize}
  \item Primero hay que preparar la versión base sobre la que luego se harán los cambios. Esa preparación se hace normalmente cuando se tiene una versión razonablemente estable, y con la que luego se quiere comparar (por ejemplo cuando le entregas a tu tutor tu primer borrador del documento completo). La preparación es sencilla: basta hacer un \texttt{make snapshot}. Eso generará un fichero \texttt{book-flatten-snapshot.tex} que contiene el estado actual del documento, expandido. Lo suyo sería luego registrarlo en el repositorio con un \texttt{cvs commit -m ``New flattened version'' book-flatten-snapshot.tex} o \texttt{git add book-flatten-snapshoot.tex; git commit -m ``New flattened version''} en git.
  \item A partir de ahí ya se puede trabajar en los cambios al documento, los que sean necesarios.
  \item Cuando se quiera obtener el fichero pdf con el control de cambios, bastará hacer un \texttt{make latexdiff}, que acabará generando \texttt{book-flatten-diff.pdf}, que será lo que estábamos buscando.
  
  \item Cuando queráis ``resetear'' el documento con el que se comparará, basta hacer un nuevo \texttt{make snapshot}.
\end{itemize}

Cristina Losada ha conseguido ejecutar el latexdiff en un windows sin necesidad de disponer de \texttt{make}, para lo cual es necesario:

\begin{itemize}
  \item Instalar un intérprete de comandos (ella usó \url{http://strawberryperl.com/releases.html}).
  
  \item MikTeX ya tiene un ejecutable de \texttt{latexdiff}.
  
  \item Guardar ``a mano'' una versión del documento correspondiente con la que compararemos los cambios futuros. Si estamos con el documento de la memoria \texttt{book.tex}, bastaría copiarlo a, por ejemplo, \texttt{book-snapshot.tex}.
  
  \item Hacer las modificaciones que procedan en el \texttt{book.tex}.
  
  \item Ejecutar en línea de comandos \texttt{latexdiff --flatten book-snapshot.tex book.tex > book-flatten-diff.tex}
  
  \item Compilar ese \texttt{book-flatten-diff.tex} que generará el pdf.
  
\end{itemize}

\subsection{Problemas conocidos}
\label{sec:problemas-conocidos}

Resumimos a continuación los problemas con los que nos hemos ido encontrando tras el uso más generalizado de esta plantilla, y, en su caso, la solución propuesta/adoptada:

\begin{itemize}

  \item Al menos en la versión 12.04 de ubuntu, se producía un fallo de compilación por un problema del \textit{option clash} del paquete \texttt{xcolor}. La solución fue incluir las opciones \texttt{[RGB,rgb]} de dicho paquete en el \texttt{documentclass}, y eliminar la inclusión de \texttt{xcolor} (que ya lo incluye, al menos, \texttt{listings}). No es muy bonito como solución, pero funcionaba. Eso dio lugar a otro problema, que hacía que todas las página aparecieran como en color cuando se llevaba a la imprenta (con el consiguiente incremento de precio). Para intentar solucionarlo, he vuelto a eliminar las opciones \texttt{[RGB,rgb]} del \texttt{documentclass}, y se las paso con un \texttt{PassOptionsToPackage}, pero está pendiente de verificación. Please, confirmadme que funciona (tanto lo del color como la compilación en un ubuntu 12.04 (o posterior).

  \item También hemos observado en la versión 12.04 de ubuntu que \texttt{evince} no es capaz de visualizar la primera página de los TFGs (generada con \texttt{tikz}), y que \texttt{xpdf} genera un core cuando intenta abrir el fichero compilado. La solución es usar \texttt{qpdfview} o \texttt{acroread}.

\end{itemize}

Si das con más problemas, escribe por favor a \contactauthor{} contándonoslos y trataremos de solucionarlo (y si tenéis la solución, contádnosla también).


\section{Ejemplos de elementos de utilidad}
\label{sec:ejempl-de-elem}

\subsection{Uso de comandos definidos}
\label{sec:uso-de-comandos}

A modo de ejemplo, hemos definido el comando \texttt{\textbackslash{}texten\{\}} en \texttt{Config/myconfig.tex} para usarlo, por ejemplo, para marcar palabras escritas en inglés (aka \texten{English}). Sigue el ejemplo para definir aquellos que utilices con frecuencia.

Si quieres escribir el símbolo \texttt{backslash} puedes usar el comando \texttt{\textbackslash{}backlash\{\}}:~\textbackslash{}.

Lo mismo aplica para el símbolo \texttt{tilde}, para lo que puedes usar el comando \texttt{\textbackslash{}textasciitilde\{\}}:~\textasciitilde{}.


\subsection{Uso de ``frases célebres''}
\label{sec:uso-de-frases}

Respecto a la frase célebre del inicio de los capítulos: todas las que hemos usado y las definiciones que las generen están sacadas del excelente trabajo de Marco Antonio Gomez-Martín y Pedro Pablo Gomez-Martín en el proyecto \texis, una plantilla para la creación de tesis y otros documentos y disponible en \cite{texis}.


\subsection{Inclusión de diagramas}
\label{sec:diagrama}

Para incluir gráficos, la compilación que utilizamos permite usar ficheros \texttt{png}, \texttt{jpg} y \texttt{pdf}, en el comando \texttt{\textbackslash{}includegraphics}. Si queréis ahorraros incluir el path a cada fichero, podéis definir todos aquellos en los que haya ficheros gráficos en el \texttt{\textbackslash{}graphicspath} del \texttt{book.tex}.

En la figura \ref{fig:fig_clobj} se muestra un ejemplo de gráfico generado automáticamente a partir de un fichero \texttt{.dia}\footnote{Tomadas de los proyectos fin de carrera de David Casillas y Manuel Villaverde.}: \texttt{diagrams/Esquema\_objetos.dia} (podéis generalizar su generación en el \texttt{Makefile}).

\begin{figure}[tphb]
  \centering \includegraphics{Esquema_objetos}
  \caption{Clasificación de los objetos para la gramática (aquí también se pueden poner acrónimos como \acs{ETTS} y símbolos como \ac{xidet})}
  \label{fig:fig_clobj}
\end{figure}


\subsection{Definición y uso de acrónimos (aquí también se pueden poner acrónimos como \acs{ETTS})}
\label{sec:uso-de-acronimos}

El uso del paquete \texttt{glossaries} permite definir los acrónimos y el sistema automáticamente gestiona su inclusión completa la primera vez que se usa. Los acrónimos de ejemplo están en el fichero \texttt{acronyms/defacronymsgl.tex} (con opciones adicionales de configuración en \texttt{acronyms/acronymsgl.tex}).

Así, si nos referimos a \ac{ETTS} o bien a \ac{EMODB}, veremos como aparecen expandidas la primera vez. A partir de ahí, sólo se usará el acrónimo como puede verse al volver a hablar de \ac{ETTS} y \ac{EMODB}.

Tiene también soporte para resetear todos los acrónimos como si no estuvieran usados. Vuelvo a incluir el párrafo anterior tras un reset (que se hace con un \texttt{\textbackslash{}glsresetall[acronym]}):

\glsresetall[acronym]

El uso del paquete acronym permite definir los acrónimos y el sistema automáticamente gestiona su inclusión completa la primera vez que se usa. Así, si nos referimos a \ac{ETTS} o bien a \ac{EMODB}, veremos como aparecen expandidas de nuevo (como si fuera la primera vez que se usan). A partir de ahí, sólo se usará el acrónimo como puede verse al volver a hablar de \ac{ETTS} y \ac{EMODB}.

Y permite también forzar que se vuelva a citar completo aunque ya se haya utilizado (con el acrónimo entre paréntesis), como puede verse en \acl{ETTS} (equivalente a \glsdesc{ETTS} que vale para cualquier glosario), y también a usar forzosamente el acrónimo. Primero reseteamos de nuevo.

\glsresetall[acronym]

Y ahora forzamos el acrónimo: \acs{EMODB} (equivalente a \glsname{EMODB} que vale para cualquier glosario). También podemos forzar a que lo ponga todo, con \acf{EMODB}.


Podemos seguir definiendo entradas de acrónimos, referirnos a \ac{DBN} por primera vez, y las siguientes aparecerá como \ac{DBN}. Pongo ahora el resto de acrónimos \ac{SQ}, \ac{EIR}, \ac{SIR} y \ac{ES}. Finalmente los repito para que se vea el efecto: \ac{SQ}, \ac{EIR}, \ac{SIR} y \ac{ES}.

Y gestiona bien los plurales, ponemos el plural como \acp{SOC} la primera vez, y luego la segunda como \acp{SOC}. Y podemos volver al singular con \ac{SOC}.


\subsection{Definición y uso de símbolos (aquí también se pueden símbolos como \ac{xidet})}
\label{sec:simbolos}

Los símbolos definidos están incluidos en el fichero \texttt{symbols/defsymbolsgl.tex} (con configuración adicional en \texttt{symbols/symbolsgl.tex}) y en esta sección mostramos algunos ejemplos.

El \ac{angstrom} se usa en biología estructural, mientras que el \ac{ohm} se usa en electrónica. También podemos poner~\ac{xdet}. Y


\section{Ejemplos de inclusión de enlaces a vídeos y listado}
\label{sec:videolink}

En marzo de 2023 recibimos una pregunta sobre la posibilidad de incluir enlaces a vídeos en el documento, y que se generara automáticamente un listado similar al de tablas, figuras o algoritmos. El código para implementarlo está también en el fichero \texttt{cover/extralistings.tex} y puedes editarlo como te interese. Para usarlo, basta con que incluyas elementos similares a \videoLink{Introducción al Laboratorio de Circuitos Electrónicos (I)}{https://www.youtube.com/watch?v=U1P2ZP0f8KM} o \videoLink{Introducción al Laboratorio de Circuitos Electrónicos (V)}{https://www.youtube.com/watch?v=nPsxMqKRGx0}.

Lo importante no es tanto que sea un listado de vídeos, sino que muestra cómo se pueden integrar listados de elementos adicionales que sean a su vez enlaces y configurar la generación del índice para generar enlaces visitables desde el mismo índice (además de su indexación convencional).

%%% Local Variables:
%%% TeX-master: "../book"
%%% End:


